\section{Training Performance and Stability Analysis}

Training is performed on historical data from the Dow Jones Industrial Average 
(DJIA) constituents using daily observations (1D frequency), covering the period 
from \texttt{2014-01-06} to \texttt{2025-12-31}. All models are
trained for up to $1{,}000{,}000$ timesteps under identical environment dynamics,
reward formulation, and preprocessing pipeline, enabling direct comparison across
algorithms and policy architectures.

Training behavior is analyzed using episodic mean reward, together with two financial 
metrics: \textit{total return}, measuring cumulative portfolio growth, and 
\textit{Sharpe ratio}, measuring risk-adjusted return. Table~\ref{tab:training_results} 
summarizes the final reward values, while Figures~\ref{fig:rollout_ep_rew_mean},
\ref{fig:financial_total_return}, and \ref{fig:financial_sharpe_ratio} illustrate the
training evolution of all monitored metrics.

\begin{figure}
    \centering
    \includegraphics[width=1\linewidth]{figures/rollout_ep_rew_mean.png}
    \caption{Training curves of episodic mean reward (\texttt{ep\_rew\_mean}) across all
    evaluated DRL algorithms and policy architectures.}
    \label{fig:rollout_ep_rew_mean}
\end{figure}

\begin{figure}
    \centering
    \includegraphics[width=1\linewidth]{figures/financial_total_return.png}
    \caption{Episode total return across all evaluated DRL algorithms and policy architectures.}
    \label{fig:financial_total_return}
\end{figure}

\begin{figure}
    \centering
    \includegraphics[width=1\linewidth]{figures/financial_sharpe_ratio.png}
    \caption{Sharpe ratio across all evaluated DRL algorithms and policy architectures.}
    \label{fig:financial_sharpe_ratio}
\end{figure}

\begin{table}[h]
\centering
\begin{tabular}{lccc}
\hline
\textbf{Model} & \textbf{Smoothed Value} & \textbf{Final Value} & \textbf{Timesteps} \\
\hline
A2C (MLP) & 0.5743 & 0.5788 & 998,400 \\
A2C (Transformer) & 0.7356 & 0.7319 & 998,400 \\
DDPG (MLP) & 2.5110 & 2.5246 & 989,280 \\
DDPG (Transformer) & 1.2787 & 1.2076 & 148,392 \\
PPO (MLP) & 1.7632 & 1.7628 & 1,001,472 \\
PPO (MLP--Transformer) & 1.7386 & 1.7396 & 1,001,472 \\
PPO (Transformer) & 2.1647 & 2.1627 & 1,001,472 \\
Recurrent PPO (LSTM) & 2.3349 & 2.3360 & 1,000,448 \\
SAC (MLP) & 1.8009 & 1.8467 & 989,280 \\
SAC (Transformer) & 0.5139 & 0.4212 & 115,416 \\
TD3 (MLP) & 2.1192 & 2.0805 & 989,280 \\
TD3 (Transformer) & 1.5086 & 1.5557 & 164,880 \\
\hline
\end{tabular}
\caption{Final training performance across different DRL algorithms and architectures.}
\label{tab:training_results}
\end{table}

\paragraph{A2C.}
Both A2C variants remain nearly flat throughout training in
Figure~\ref{fig:rollout_ep_rew_mean}, converging to relatively low reward levels.
This behavior is consistently reflected in total return and Sharpe ratio
(Figures~\ref{fig:financial_total_return} and
\ref{fig:financial_sharpe_ratio}), indicating limited policy improvement during
training. The Transformer variant provides only a marginal performance gain over the
MLP baseline.

\paragraph{PPO.}
All PPO variants exhibit stable learning dynamics with clear convergence behavior across
all monitored metrics. The Transformer policy achieves the strongest performance within
the PPO family (Table~\ref{tab:training_results}), while also maintaining strong total
return and Sharpe ratio behavior. The hybrid MLP--Transformer architecture performs
similarly to the MLP baseline, suggesting limited benefit from partial separation of
portfolio and market representations in this experimental setting.

\paragraph{Recurrent PPO (LSTM).}
Recurrent PPO achieves the highest episodic reward
(Table~\ref{tab:training_results}); however, the reward trajectory is non-monotonic and
exhibits noticeable oscillations
(Figure~\ref{fig:rollout_ep_rew_mean}), indicating reduced training stability compared
to standard PPO variants. This behavior is further reflected in the total return metric,
where extreme positive spikes approaching numerical divergence appear
(Figure~\ref{fig:financial_total_return}), suggesting episodic instability under
long-horizon portfolio compounding.

\paragraph{DDPG.}
DDPG (MLP) achieves the strongest overall performance across all monitored metrics.
It consistently attains the highest total return while maintaining strong Sharpe ratio
behavior and stable reward progression. In contrast, the Transformer variant terminates
early at approximately 148k timesteps and performs substantially worse across all
metrics.

\paragraph{SAC.}
The SAC MLP model demonstrates moderate but relatively stable performance across reward,
total return, and Sharpe ratio. The Transformer variant fails to complete training,
terminating at approximately 115k timesteps with weaker performance and unstable
learning behavior.

\paragraph{TD3.}
The TD3 MLP model exhibits non-monotonic reward dynamics, starting from relatively high
reward values before gradually declining during training. The Transformer variant
terminates early at approximately 165k timesteps and additionally exhibits extreme
spikes in total return
(Figure~\ref{fig:financial_total_return}), similar to recurrent PPO, indicating
episodic instability rather than sustained portfolio growth.

Overall, the results show consistent trends across reward and financial metrics.
MLP policies provide the most stable training behavior, while Transformer architectures
improve representation capacity for on-policy methods---particularly PPO---but reduce
stability in off-policy algorithms. Recurrent policies provide strong temporal modeling
capabilities, although they introduce additional instability in long-horizon portfolio
dynamics.